\chapter*{Glossary}
\addcontentsline{toc}{chapter}{Glossary}
\markboth{Glossary}{Glossary}

\noindent The terms below appear in this booklet. Definitions are given
in the sense used here; broader uses may vary.

\bigskip

\noindent\textbf{Data.}
A symbol, signal, or fact placed within a framework of meaning, purpose,
or use, relative to which it is relevant and interpretable. Without such
a framework it is noise. See also \textit{DIKW pyramid}.

\medskip

\noindent\textbf{Datum.}
The Latin singular of \textit{data} (``something given''). Past participle
of \textit{dare}, to give.

\medskip

\noindent\textbf{DIKW pyramid.}
The hierarchy Data -- Information -- Knowledge -- Wisdom. Proposed by
Ackoff (1989) as a way of naming the increasing levels of organisation
and meaning that can be applied to raw symbols.

\medskip

\noindent\textbf{Episodic memory.}
The system that encodes and retrieves personally experienced events,
supported by the medial temporal lobe (MTL). Works alongside the
hippocampus to convert short-term experience into retrievable
long-term memories.

\medskip

\noindent\textbf{Hippocampus.}
A small brain structure central to memory consolidation --- the
``warehouse manager'' of the analogy in Chapter~2. Converts
short-term material into long-term storage and organises retrieval.

\medskip

\noindent\textbf{Information.}
Raw data that has been extracted, organised, and placed in a context
that gives it meaning within a defined structure (e.g., ``the February
Revolution of 1848'' rather than a bare year and month).

\medskip

\noindent\textbf{Knowledge.}
The understanding that arises when information is combined with the
relationships among its parts, yielding a grasp of the \emph{why}
and \emph{how}, not merely a list of facts.

\medskip

\noindent\textbf{Mental representation.}
An intentional mental state --- thought, belief, perception, or
imagining --- that is ``about'' something in the world and can be
evaluated for truth, accuracy, or appropriateness.

\medskip

\noindent\textbf{Representation.}
One thing standing in the place of another: a word for an object,
a number for a quantity, a name for a historical event. The four
criteria (connection, causal role, exclusivity, teleological
justification) define what makes something a genuine representation
rather than a mere correlation.

\medskip

\noindent\textbf{Sign.}
Any mark, word, or signal used to stand for something other than itself.
A name like ``the July Monarchy'' is a sign: it has no inherent
connection to July, but everyone accepts it as pointing to a specific
set of historical events.

\medskip

\noindent\textbf{Wisdom.}
The highest layer of the DIKW pyramid: the capacity to apply
knowledge in the form of sound judgment and correct action, rather
than merely possessing the underlying facts or understanding.
